\documentclass[a4paper]{article}

%% Language and font encodings
\usepackage[portuges]{babel}
\usepackage{fontspec}

%% Sets page size and margins
\usepackage[a4paper,top=3cm,bottom=2cm,left=3cm,right=3cm,marginparwidth=1.75cm]{geometry}

%% Useful packages
\usepackage{amsmath,amsthm,amssymb,amsfonts}
\usepackage{graphicx}
\usepackage[colorinlistoftodos]{todonotes}
\usepackage[colorlinks=true, allcolors=blue]{hyperref}

\newcommand{\R}{\mathbb{R}}
\newcommand{\N}{\mathbb{N}}
\newcommand{\Z}{\mathbb{Z}}
\providecommand{\C}{\mathbb{C}}

\theoremstyle{definition}
\newtheorem{defin}{Definição}

\theoremstyle{plain}
\newtheorem{theorem}[defin]{Teorema}
\newtheorem{corollary}[defin]{Corolário}



\title{Laboratório de Matemática --- Ficha N}

\author{Número: xxxxx Nome\\
        Número: yyyyy Nome\\
}

\date{dd/mm/yyyy}

\begin{document}
\maketitle

%\begin{abstract}
%Your abstract.
%\end{abstract}

\section{Introdução do tema}

\emph{\small Incluir aqui uma breve descrição do tema}

\begin{defin}
Um triângulo diz-se \emph{rectângulo} se tiver um ângulo recto.
\end{defin}

\begin{theorem}[Pitágoras]
Num triângulo rectângulo de lados $a$, $b$ e $c$, em que $c$ é o lado maior, tem-se a igualdade
$$ a^2 + b^2 = c^2 $$
\end{theorem}

\begin{corollary}
No plano cartesiano $\R^2$, a distância entre os pontos $A=(a_1,a_2)$ e $B=(b_1,b_2)$ é igual a
$$ \sqrt{(a_1-b_1)^2 + (a_2-b_2)^2} $$
\end{corollary}
\section{Trabalho realizado}


%Podem inserir gráficos, usando \includegraphics: por exemplo,
%\includegraphics[width=0.5\linewidth]{cloud}
%
% Veremos na aula como fazer upload de imagens para usar aqui

Usando o Geogebra, realizámos os seguintes passos:

\begin{enumerate}
\item usando a linha de \emph{input}, construímos dois pontos \texttt{A=(a1,a2)} e \texttt{B=(b1,b2)}
\item construímos um ponto $C$, de coordenadas $(b1,a2)$
\item observamos que o triângulo $[ABC]$ é rectângulo em $C$
\item calculámos a distância entre $A$ e $B$, usando o comando \texttt{d1 = Distance(A,B)}.
\item calculámos a expressão $\sqrt{(a_1-b_1)^2 + (a_2-b_2)^2}$, usando o commando\\ \texttt{d2=sqrt((x(A)-x(B))\^2+sqrt(y(A)-y(B))\^2}

\item fazendo variar os pontos $A$ e $B$ verificámos que os valores de \texttt{d1} e de \texttt{d2} são sempre iguais entre si.
\end{enumerate}

\section{Conclusões}

As nossas experiências são consistentes com a fórmula conhecida para a distância entre dois pontos.

%expor aqui as conclusões que retiraram do trabalho

\end{document}
